% ══════════════════════════════════════════════════════════════════
% CHAPITRE 1 — TRAVAUX EXISTANTS
% Plan : Intro · 1.1 · 1.2 · 1.3 · 1.4 · Conclusion
% ══════════════════════════════════════════════════════════════════

\pagechapitre{chap:travaux}{Travaux existants}

\section*{Introduction}
\addcontentsline{toc}{section}{Introduction}

Ce chapitre dresse un état des lieux des connaissances et des
technologies sur lesquelles s'appuie notre travail. Il présente
d'abord l'agriculture de précision et le suivi parcellaire, puis
la nature des données agricoles exploitables. Il expose ensuite
les principales techniques d'Intelligence Artificielle appliquées
à l'agriculture, avant d'analyser les solutions existantes afin
de positionner notre contribution.

% ══════════════════════════════════════════════════════════════════
\section{Agriculture de précision et suivi parcellaire}
\label{sec:agri_precision}
% ══════════════════════════════════════════════════════════════════

\subsection{Définition et principes}

L'agriculture de précision désigne un ensemble de pratiques et de
technologies visant à optimiser la gestion des cultures en tenant
compte de la variabilité spatiale et temporelle des parcelles
\cite{liakos2018machine}. Plutôt que d'appliquer uniformément les
intrants (eau, engrais, produits phytosanitaires), elle propose
d'ajuster les interventions aux besoins réels de chaque zone,
mesurés à partir de données objectives. Ses bénéfices attendus
sont triples : augmentation des rendements, réduction des coûts
et préservation de l'environnement.

\subsection{Le suivi parcellaire}

Le suivi parcellaire consiste à observer régulièrement l'état
d'une parcelle agricole --- humidité du sol, état sanitaire des
plantes, stade de développement --- afin de déclencher au bon
moment les interventions appropriées : irrigation, fertilisation
ou traitement. Traditionnellement assuré par l'observation
visuelle de l'agriculteur et les conseils d'agronomes, ce suivi
peut aujourd'hui être outillé par des capteurs, des stations
météorologiques, des images satellitaires et des applications
d'aide à la décision.

\subsection{Contexte agricole du Tchad}

Le Tchad dispose d'environ 39 millions d'hectares de terres
cultivables, dont une faible fraction seulement est exploitée
\cite{banquemondiale2023}. L'agriculture y est essentiellement
pluviale et familiale, structurée autour de deux grandes zones
agro-climatiques : la zone \textbf{sahélienne} au centre
(N'Djamena, Abéché), caractérisée par une pluviométrie faible et
irrégulière (300--600 mm/an), et la zone \textbf{soudanienne} au
sud (Sarh, Moundou, Bongor), plus arrosée (800--1~200 mm/an).
Les cinq cultures retenues dans ce travail --- mil, sorgho,
arachide, maïs et coton --- représentent l'essentiel des
superficies cultivées du pays.

\subsection{Défis du suivi parcellaire en zone sahélienne}

Le déploiement du suivi parcellaire au Tchad se heurte à
plusieurs obstacles :

\begin{itemize}[leftmargin=1.2cm]
	\item le nombre limité d'agronomes et de techniciens agricoles
	en zone rurale ;
	\item l'absence de stations météorologiques de proximité et
	d'accès aux prévisions localisées ;
	\item le coût élevé des équipements de mesure (sondes
	d'humidité, capteurs NPK) ;
	\item la faible couverture internet et électrique des zones
	agricoles.
\end{itemize}

Ces contraintes orientent naturellement vers des solutions
logicielles légères, exploitant des données gratuites
(satellites, réanalyses climatiques) et accessibles depuis un
simple smartphone.

% ══════════════════════════════════════════════════════════════════
\section{Données agricoles : nature et hétérogénéité}
\label{sec:donnees_agricoles}
% ══════════════════════════════════════════════════════════════════

Les données exploitables pour le suivi parcellaire sont de
natures très diverses ; leur hétérogénéité constitue à la fois
une richesse et un défi pour les systèmes d'aide à la décision.

\subsection{Données climatiques}

Les variables climatiques --- température, humidité relative,
précipitations, rayonnement solaire, vitesse du vent et
évapotranspiration potentielle (ETP) --- conditionnent
directement les besoins en eau des cultures. Deux sources
gratuites sont particulièrement pertinentes pour le Tchad :

\begin{itemize}[leftmargin=1.2cm]
	\item \textbf{NASA POWER} \cite{stackhouse2019nasa} : base de
	réanalyses climatiques mondiales fournissant des séries
	journalières historiques depuis 1981, à une résolution
	de 0,5° × 0,5°, utilisée ici pour constituer le jeu
	d'entraînement du module Irrigation ;
	\item \textbf{Open-Meteo} \cite{zippenfenig2023openmeteo} :
	API météorologique temps réel et gratuite, sans clé
	d'authentification, utilisée par l'application pour
	pré-remplir automatiquement les formulaires.
\end{itemize}

\subsection{Données pédologiques}

Les propriétés du sol --- teneurs en azote (N), phosphore (P) et
potassium (K), pH, texture (sableux, limoneux, argileux) et
humidité --- déterminent la disponibilité des nutriments et le
choix des engrais. Ces données proviennent d'analyses de sol en
laboratoire ou de kits de terrain.

\subsection{Données visuelles et télédétection}

L'état sanitaire des cultures peut être évalué à partir d'images :

\begin{itemize}[leftmargin=1.2cm]
	\item \textbf{Photographies de feuilles} prises au smartphone,
	analysables par des réseaux de neurones convolutifs pour
	détecter les symptômes de maladies
	\cite{mohanty2016deep} ;
	\item \textbf{Images satellitaires} : le programme européen
	Sentinel-2 \cite{sentinel2esa} fournit gratuitement des
	images multispectrales à 10~m de résolution, revisitées
	tous les 5 jours, dont on dérive des indices de
	végétation ;
	\item \textbf{Imagerie aérienne} par drones multispectraux,
	offrant une résolution centimétrique, mais dont le coût
	d'acquisition reste élevé pour le contexte tchadien.
\end{itemize}

\subsection{L'indice de végétation NDVI}

Le NDVI (\textit{Normalized Difference Vegetation Index})
\cite{rouse1974monitoring} est l'indice le plus utilisé pour
quantifier la vigueur de la végétation. Il est calculé à partir
des réflectances dans le rouge (RED) et le proche infrarouge
(NIR) :

\begin{equation}
	\text{NDVI} = \frac{\text{NIR} - \text{RED}}
	{\text{NIR} + \text{RED}}
	\label{eq:ndvi}
\end{equation}

Ses valeurs varient de $-1$ à $+1$ : une végétation dense et
saine réfléchit fortement le proche infrarouge (NDVI $> 0{,}5$),
tandis qu'une végétation stressée ou malade voit son NDVI chuter
(NDVI $< 0{,}3$).

% ══════════════════════════════════════════════════════════════════
\section{Intelligence artificielle appliquée à l'agriculture}
\label{sec:ia_agriculture}
% ══════════════════════════════════════════════════════════════════

\subsection{Apprentissage automatique classique}

Les algorithmes d'apprentissage supervisé classiques excellent
sur les données tabulaires (climatiques, pédologiques). Deux
familles sont particulièrement adaptées à notre problème :

\begin{itemize}[leftmargin=1.2cm]
	\item \textbf{Random Forest} \cite{breiman2001random} :
	ensemble d'arbres de décision entraînés sur des
	sous-échantillons aléatoires, dont la prédiction finale
	résulte d'un vote majoritaire. Robuste au bruit et peu
	sensible au surapprentissage, il fournit en outre une
	mesure d'importance des variables ;
	\item \textbf{XGBoost} \cite{chen2016xgboost} : implémentation
	optimisée du \textit{gradient boosting}, construisant
	séquentiellement des arbres correcteurs d'erreurs.
	Il domine régulièrement les compétitions de science des
	données sur données structurées.
\end{itemize}

\subsection{Apprentissage profond et vision par ordinateur}

Les réseaux de neurones convolutifs (CNN) ont révolutionné la
reconnaissance d'images. Appliqués à la santé des plantes,
Mohanty \textit{et al.} \cite{mohanty2016deep} ont atteint
99,35\,\% de précision sur le dataset PlantVillage (54~000
images en conditions contrôlées) ; toutefois, la précision de
leur modèle chute à environ 31\,\% sur des images de terrain,
illustrant le fossé entre conditions de laboratoire et réalité
agricole --- un point central pour notre travail sur des images
tchadiennes réelles.

\subsection{Transfert d'apprentissage et modèles légers}

Le \textit{transfer learning} consiste à réutiliser un réseau
pré-entraîné sur un grand corpus générique (ImageNet, 14
millions d'images) et à ne réentraîner que ses couches
supérieures sur le jeu de données cible. Cette approche est
incontournable lorsque les données annotées sont rares, comme
c'est le cas pour les cultures tchadiennes.
\textbf{MobileNetV2} \cite{sandler2018mobilenetv2} est une
architecture conçue pour les environnements contraints :
ses blocs résiduels inversés réduisent drastiquement le nombre
de paramètres, et sa conversion au format \textbf{TensorFlow
	Lite} permet un déploiement léger (quelques mégaoctets) adapté
au web et au mobile.

% ══════════════════════════════════════════════════════════════════
\section{Analyse des solutions existantes et positionnement}
\label{sec:solutions_existantes}
% ══════════════════════════════════════════════════════════════════

\subsection{Solutions existantes}

Plusieurs applications d'agriculture numérique sont déployées en
Afrique et dans le monde. Le tableau~\ref{tab:solutions} en
compare les principales caractéristiques avec notre proposition.

\begin{table}[H]
	\centering
	\caption{Comparaison des solutions existantes avec Agro-IA Tchad}
	\label{tab:solutions}
	\small
	\begin{tabularx}{\textwidth}{lXXXX}
		\toprule
		\textbf{Critère} & \textbf{Plantix} & \textbf{Esoko} &
		\textbf{PlantVillage Nuru} & \textbf{Agro-IA Tchad} \\
		\midrule
		Origine & Allemagne & Ghana & International & Tchad \\
		Détection maladies & Oui (photo) & Non & Oui (photo) &
		Oui (photo + NDVI) \\
		Conseil irrigation & Non & Non & Non & Oui (IA) \\
		Conseil fertilisation & Partiel & Non & Non & Oui (IA) \\
		Historique local & Non & Non & Non & Oui (SQLite) \\
		Cultures tchadiennes & Non & Non & Partiel & Oui (5 cultures) \\
		Météo temps réel & Oui & Oui & Non & Oui (Open-Meteo) \\
		Coût & Gratuit & Abonnement & Gratuit & Gratuit \\
		\bottomrule
	\end{tabularx}
\end{table}

\subsection{Limites des solutions existantes}

L'analyse du tableau~\ref{tab:solutions} révèle trois lacunes
majeures pour le contexte tchadien : (i) aucune solution ne
couvre simultanément l'irrigation, la fertilisation et la
détection des maladies ; (ii) les modèles de reconnaissance sont
entraînés sur des cultures et des conditions étrangères, peu
représentatives des variétés et symptômes locaux ; (iii) aucune
ne propose de traçabilité locale des analyses permettant un
suivi historique des parcelles.

\subsection{Positionnement de notre contribution}

Notre système Agro-IA Tchad se positionne comme une réponse
intégrée à ces lacunes :

\begin{itemize}[leftmargin=1.2cm]
	\item \textbf{Multi-modules} : quatre modules complémentaires
	(Irrigation, Fertilisation, Détection des maladies,
	Historique) réunis dans une même application ;
	\item \textbf{Adaptation locale} : modèles entraînés sur des
	données climatiques tchadiennes (NASA POWER, 5 villes)
	et sur un dataset d'images collectées au Tchad
	(287 images, 10 classes) ;
	\item \textbf{Traçabilité} : base SQLite intégrée avec
	statistiques et export CSV ;
	\item \textbf{Accessibilité} : application web gratuite,
	déployée sur le cloud, exploitant la météo temps réel.
\end{itemize}

\section*{Conclusion}
\addcontentsline{toc}{section}{Conclusion}

Ce chapitre a présenté les fondements de l'agriculture de
précision, la diversité des données agricoles mobilisables, les
techniques d'IA pertinentes --- Random Forest, XGBoost et CNN
par transfert d'apprentissage --- ainsi que les limites des
solutions existantes au regard du contexte tchadien. Ce
positionnement établi, le chapitre suivant procède à l'analyse
des besoins et à la conception détaillée du système Agro-IA
Tchad.